\documentclass[]{report}

\usepackage{import}

\import{../}{settings.tex}

\title{Cours sur les séries pour l'agrégation}

\author{Cours de Stephane VASSOULT, transcrit par SADR TAHOURI Luca}
\date{}

\begin{document}

\maketitle
\tableofcontents

\chapter{Séries Numériques}

\section{Généralités}

\begin{df}{Série numérique, somme partielle, reste d'une série}{}
    On appelle série numérique une suite dans $\K$ ($=\R$ ou $\C$) de la forme
    \[\fundf{\N}{\K}{n}{S_n=\sum_{k=0}^{n}u_k}\]
    Où $u$ est une autre suite numérique. On l'appelle série de terme général $u_n$ et $S_n$ la somme partielle de la série. Si la série converge vers $S$, on appelle somme de la série la limite. Dans ce cas, on désigne par
    \[R_n=S-S_n=\sum_{k=n+1}^{+\infty}u_k\]
    le reste de la série (C'est une suite qui tend vers $0$).
\end{df}

\begin{re}{}{}
    \begin{itemize}
        \item Les séries sont des suites particulières
        \item Toute suite peut s'écrire comme une série.
    \end{itemize}
    Ce qui diffère: Les moyens d'études
\end{re}

\begin{prop}{Conditions de convergence des séries numériques}{}
    Une série numérique (Dans $\R$ ou $\C$) est convergente si la suite des sommes partielles $(S_n)_\nN$ est une suite de Cauchy.
\end{prop}

\begin{cor}{}{}
    Première conséquence: si la série converge, alors son terme général tend vers $0$
\end{cor}

\begin{df}{Série qui diverge grossièrement}{}
    Si le terme général de la série ne tend pas vers $0$, on dit que la série diverge grossièrement.
\end{df}

\section{Séries à termes positifs}

\begin{df}{Série à termes positifs}{}
    Une série à termes positifs est une série $\sum u_n$ avec $u_n \ge 0$. Sa suite $(S_n)_\nN$ est croissante donc
    \begin{itemize}
        \item Si $(S_n)_\nN$ est majorée, elle converge
        \item Sinon $S_n\toninfty +\infty$
    \end{itemize}
\end{df}

\begin{prop}{Critère de comparaison des séries numériques}{}
    Soient $\sum u_n$ et $\sum v_n$ deux séries à termes positifs. Si $u_n\le v_n$ à partir d'un certain rang, alors
    \begin{itemize}
        \item Si $\sum u_n$ diverge, $\sum v_n$ diverge
        \item Si $\sum v_n$ converge, alors $\sum u_n$ converge
    \end{itemize}
    Séries avec lesquelles conparer:
    \begin{itemize}
        \item Séries Géométriques: $a\ge 0$
              \[\sum a^n \cho{\text{ converge si }a<1\\ \text{ diverge si }a \ge 1}\]
        \item Séries de Riemann:
              \[\sum \frac{1}{n^{\alpha}}\cho{\text{ converge si }\alpha>1\\ \text{ diverge si }\alpha < 1}\]
        \item Séries de Bertrand
              \[\sum \frac{1}{n^\alpha (\ln (n))^\beta}\cho{\text{converge si }\alpha>1 \text{ ou }(\alpha=1 \text{ et } \beta >1)\\ \text{diverge si $a<1$, $a=1$ et $\beta \le 1$}}\]
    \end{itemize}
    VOIR PHOTO POUR COMPARAISON
\end{prop}

\begin{cor}{Déclinaison des critères de comparaison}{}
    \begin{itemize}
        \item Si $u_n \sim v_n$, alors $\sum u_n$ et $\sum v_n$ sont de même nature.
        \item Si $u_n= o(v_n)$ ou $u_n=O(v_n)$ et $\sum v_n$ converge, alors $\sum u_n$ converge.
        \item Critère de Cauchy (ou régle de Cauchy):
              \[\sqrt[n]{u_n}\to l\implies \cho{\text{Si $l>1$, alors $\sum u_n$ diverge}\\ \text{Si $l<1$ alors $\sum u_n$ converge} \\ \text{Si $l=1$, pas de conclusion}}\]
        \item Critère de d'Alembert (ou règle de d'Alembert): Si $u_n \neq 0$ APCR
              \[\frac{u_{n+1}}{u_n}\toninfty l\implies \cho{\text{Si $l>1$, alors $\sum u_n$ diverge}\\ \text{Si $l<1$ alors $\sum u_n$ converge} \\ \text{Si $l=1$, pas de conclusion}}\]
    \end{itemize}
\end{cor}

\section{Séries à terme de signe quelconque}

A priori, $(u_n)_\nN$ est une suite réelle ou complexe quelconque.

\begin{prop}{Absolue convergence}{}
    $\sum |u_n|$ converge alors $\sum u_n $ converge
\end{prop}

Si on ne peut pas conclure, on peut utiliser d'autres outils:
\begin{itemize}
    \item Séries alternées
    \item Transformation d'Abel
    \item Développement asymptotique
\end{itemize}

\begin{dfprop}{Série alternée}{}
    $\sum u_n$ est une série alternée si
    \begin{itemize}
        \item $\exists (a_n)_\nN$ une suite positive décroissante qui tend vers $0$ tq $u_n=(-1)^n a_n$. Dans ce cas la série de terme général $u_n$ est convergente.
    \end{itemize}
\end{dfprop}

\begin{dfprop}{Transformation d'Abel}{}
    $(u_n)_\nN$ et $(v_n)_\nN$ deux suites. On note $U_n=\sum_{k=0}^{n}$ et $V_n=\sum_{k=0}^{n} v_k$ alors pour tout $p,q$ entiers, $p\ge q$,
    \[\sum_{n=q+1}^{p} u_n v_n=\sum_{n=q+1}^{p-1}(u_n-v_{n+1}) V_n +u_p V_p -u_{q+1} V_q\]
    Application: Si $(V_n)$ est bornée et si $(u_n)_\nN$ est une suite positive décroissante, qui tend vers $0$, alors $\sum u_nv_n$ converge.
\end{dfprop}

\begin{dfprop}{Développement asymptotique}{}
    Si $u_n=v_n+o (w_n)$ avec $\sum w_n$ une série absolument convergente, alors $\sum u_n$ et $\sum v_n$ sont de même nature.
\end{dfprop}

\subsection*{Sommes partielles, restes}

Pour les séries à termes positifs:\\
On suppose que $u_n\ge 0$ et $\sum u_n$ diverge\\
\[u_n\sim v_n\implies  \sum_{k=0}^{n} v_k \sim \sum_{k=0}^{n} u_k\]
\[u_n=o(v_n)\implies \sum_{k=0}^{n} u_k = o\left( \sum_{k=0}^{n} v_k\right)\]
On suppose que $v_n\ge 0$ et $\sum v_n$ converge, alors\\
\[u_n \sim v_n\implies \sum_{k=n+1}^{+\infty}\sim \sum_{k=n+1}^{+\infty} v_k\]
\[u_n=o(v_n) \implies \sum_{k=n+1}^{+\infty}= o \left(\sum_{k=n+1}^{+\infty} v_k\right)\]
VOIR PHOTOS

\subsection*{Résultats sur les séries de Riemann et Bertrand}

Si $\funto{f}{]0,+\infty[}{\R}$ est une fonction positive décroissante et localement integrable.
\[f(k)\ge \int_{k}^{k+1} f(t)\, dt \ge f(k+1)\]
\[\sum_{k=2}^{n}\int_{k-1}^{k} f(t)\, dt \ge \sum_{k=1}^{n} f(k) \ge \sum_{k=1}^{n} \int_{k}^{k+1} \int_{k}^{k+1} f(t)\, dt\]
Donc la série $\sum f(k)$ et l'intégrale généralisée $\int_{1}^{+\infty}f$ sont de même nature. Ces inégalités fournissent aussi des équivalents

\subsection*{Produit de Cauchy de deux séries numériques}

\begin{df}{Produit de Cauchy de deux séries numériques}{}
    Si $\sum u_n,\sum v_n$ sont deux séries numériques, on appelle produit de Cauchy des deux séries la séries $\sum w_n$ de terme géénral
    \[w_n =\sum_{k=0}^{n} u_k v_{n-k}\]
\end{df}

\begin{tm}{Théorème de Mertens}{}
    Si $\sum u_n$ converge et $\sum v_n$ sont absolument convergentes, alors $\sum w_n$ est convergente et
    \[\sum_{n=0}^{+\infty} w_n=\left(\sum_{n=0}^{+\infty} u_n\right)\left(\sum_{n=0}^{+\infty} v_n\right)\]
    En général, la convergence simple de $\sum u_n$ et $\sum v_n$ ne suffit pas à assurer la convergence de $\sum w_n$.\\
    Conséquence: Le produit de Cauchy de deux séries absoluement convergente est une série absoluement convergente
\end{tm}

\begin{proof}{Théorème de Mertens}{}
    On pose $U_n=\sum_{k=0}^{n} u_k \to U$ et $V_n=\sum_{k=0}^{n} v_k\to V$ et
    \[W_n =\sum_{k=0}^{n} =\sum_{k=0}^{n} \sum_{i+j=k}U_i V_j\]
    \begin{align*}
        \sum_{i+j\le n} U_i V_j   & = \sum_{i+j\ge n} u_i \left( V_j- V_{j-1}\right)          \\
                                  & =\sum_{i+j=n} u_i V_j                                     \\
        \left| W_n - U_n V\right| & =\left|\sum_{i+j=n} u_i V_j - \sum_{i=0}^n u_i V\right|   \\
                                  & =\left| \sum_{i=0}^{n} u_i \left( V_{i-1}-V\right)\right|
    \end{align*}
\end{proof}

\begin{prop}{Rabe-Duhamel}{}
    $u_n$ une suite et $\sum v_n$ absoluement convergent.
    \[\frac{u_{n+1}}{u_n}=1-\frac{\alpha}{n} + o\left(\frac{1}{n}\right)\]
    Alors si $\alpha>1,\sum u_n$ converge, si $\alpha<1$, alors $\sum u_n $ diverge.\\
    Si $\frac{u_{n+1}}{u_n}=1-\frac{\alpha}{n}+o(v_n)$, alors $u_n \sim \frac{1}{n^\alpha}$
\end{prop}
\end{document}